\worksheettitle
  {Классы многозначных чисел}
  {\modulemark{M03} Версия учителя}

\begin{teacherbox}[Паспорт модуля]
  \textbf{Результат:} ученик делит запись числа справа налево на группы по три
  цифры, называет классы, различает группу класса и количество его единиц,
  восстанавливает число по классам.

  \textbf{Предварительно:} ученик сохраняет нулевые разряды при сборке числа.
  Время не ограничивается; решение о переходе принимается по контрольной точке.
\end{teacherbox}

\begin{checkpointbox}[Вход: решение]
  \(60\,000+3\,000+40+2=63\,042.\)

  Ошибка \(6\,342\) показывает потерю нулевого разряда. До M03 используйте
  M02-R.
\end{checkpointbox}

\section{Открытие способа}

Для числа \(7204061\) последовательно выделяем справа \(061\), затем \(204\),
затем \(7\):
\[
  7204061=7\,|\,204\,|\,061=7\,204\,061.
\]

\begin{teacherbox}[Точная формулировка]
  Запись числа делят на классы справа налево, по три цифры. Только крайняя
  левая группа может содержать одну или две цифры. Не используйте выражения
  «справа от первого класса» и «непервый класс»: I класс расположен справа.
\end{teacherbox}

Проба с опорой:
\begin{enumerate}
  \item \(5308=5\,|\,308=5\,308\);
  \item \(604015=604\,|\,015=604\,015\);
  \item \(80304072=80\,|\,304\,|\,072=80\,304\,072\).
\end{enumerate}

\section{Названия и содержание классов}

Классы справа налево: I класс единиц, II класс тысяч, III класс миллионов.

Для числа \(325\,041\,708\):
\begin{itemize}
  \item группа класса миллионов \(325\) обозначает \(325\) единиц класса
    миллионов;
  \item группа класса тысяч \(041\) обозначает \(41\) единицу класса тысяч;
  \item группа класса единиц \(708\) обозначает \(708\) единиц класса единиц.
\end{itemize}

\begin{teacherbox}[Терминологическая опора]
  «Группа класса» означает цифры в записи числа. «Количество единиц класса»
  означает число, полученное после отбрасывания начальных нулей. Например,
  группа \(006\) класса тысяч обозначает \(6\) единиц этого класса.

  Не смешивайте это с количеством всех полных тысяч в числе:
  \(325\,041\,708\) содержит \(325\,041\) полную тысячу, хотя в группе класса
  тысяч записано \(041\), то есть \(41\) единица класса тысяч.
\end{teacherbox}

\newpage
\section{Ответы к закреплению}

Таблица для \(52\,006\,090\): группы \(52\), \(006\), \(090\); количества
единиц классов \(52\), \(6\), \(90\). Удаление начальных нулей внутри записи
сдвинуло бы последующие цифры в другие разряды.

\begin{teacherbox}[Разделение на классы]
  \begin{enumerate}
    \item \(4\,072\);
    \item \(53\,008\);
    \item \(604\,015\);
    \item \(7\,030\,406\);
    \item \(80\,005\,009\);
    \item \(302\,000\,070\).
  \end{enumerate}
\end{teacherbox}

Для числа \(418\,027\,305\): III класс миллионов, группа \(418\); II класс
тысяч, группа \(027\), \(27\) единиц класса тысяч; I класс единиц, группа
\(305\).

Нулевые классы:
\begin{itemize}
  \item \(52\,000\,308\): класс тысяч;
  \item \(407\,006\,000\): класс единиц;
  \item \(9\,000\,040\): класс тысяч.
\end{itemize}

\begin{teacherbox}[Диагностика]
  Если ученик уверенно ставит пробелы, но не называет классы, вернитесь к
  нумерации от I до III. Если теряет нули, свяжите группу класса с тремя разрядами.
  Если начинает слева, предложите закрыть левую часть записи и выделить
  последние три цифры.
\end{teacherbox}

\section{Обратное действие}

\begin{enumerate}
  \item \(36\,|\,005\,|\,008=36\,005\,008\);
  \item \(4\,|\,000\,|\,072=4\,000\,072\);
  \item \(205\,|\,009\,|\,040=205\,009\,040\).
\end{enumerate}

\begin{teacherbox}[Что спросить]
  «Сколько цифр должно быть в группе тысяч?», «Чем дополнить число \(5\) до
  трёх цифр?», «Что произойдёт с разрядами, если группу \(000\) пропустить?»
\end{teacherbox}

\newpage
\section{Ошибки и сравнение}

Исправления:
\begin{enumerate}
  \item \(7204061=7\,|\,204\,|\,061\): деление начали слева;
  \item \(53008=53\,|\,008\): в правой группе оставили две цифры;
  \item \(42\) миллиона, \(7\) тысяч, \(90\) единиц:
    \(42\,|\,007\,|\,090=42\,007\,090\): группы не дополнены нулями.
\end{enumerate}

Сравнение:
\[
  8\,041\,005<8\,401\,005,
  \qquad
  27\,006\,050<27\,060\,005.
\]
В обеих парах сначала различаются группы класса тысяч.

\begin{teacherbox}[Важно]
  При сравнении классы рассматривают слева направо, а при делении записи на
  классы действуют справа налево. Проговорите различие этих двух действий.
\end{teacherbox}

\section{Контрольная точка: ответы}

\begin{enumerate}
  \item \(90305008=90\,|\,305\,|\,008=90\,305\,008\).
  \item Группы: миллионы \(90\), тысячи \(305\), единицы \(008\).
  \item В группе класса тысяч записано \(305\) единиц класса; в группе класса
    единиц --- \(8\) единиц класса.
  \item \(17\,|\,004\,|\,006=17\,004\,006\).
  \item Группы \(004\) и \(006\) занимают по три разряда и сохраняют позиции
    цифр \(4\) и \(6\).
  \item \(6150302=6\,|\,150\,|\,302=6\,150\,302\); исходную запись ошибочно
    делили слева.
\end{enumerate}

\begin{resultbox}[Решение о переходе]
  \begin{itemize}
    \item \textbf{Переход к M04:} верны пункты с 1-го по 5-й, объяснение связано с
      тремя разрядами класса.
    \item \textbf{M03-R:} две или более ошибки либо повторяется одна системная
      ошибка: начало слева, неполная группа или потеря нулей.
    \item \textbf{M03-H:} способ понятен, но требуется самостоятельное
      закрепление деления и сборки числа по классам.
    \item \textbf{M03\(\star\):} базовый результат достигнут устойчиво.
  \end{itemize}
\end{resultbox}

\newpage
\section{Маршрут после коррекции}

\begin{teacherbox}[Повторная проверка]
  Дайте новое число \(70402009\). Ожидается \(70\,|\,402\,|\,009\). Ученик
  называет группы \(70\), \(402\), \(009\), сообщает \(402\) единицы класса
  тысяч и \(9\) единиц класса единиц, затем собирает \(12\) миллионов,
  \(5\) тысяч и \(30\) единиц:
  \(12\,005\,030\).
\end{teacherbox}

\begin{teacherbox}[Критерии устойчивости]
  Ученик:
  \begin{itemize}
    \item начинает деление записи с правого края;
    \item сохраняет три цифры во всех группах, кроме крайней левой;
    \item правильно называет I, II и III классы;
    \item различает группу \(007\) и \(7\) единиц класса тысяч;
    \item выполняет обратное действие без образца;
    \item объясняет ошибку через разряды, а не только исправляет ответ.
  \end{itemize}
\end{teacherbox}

\begin{notebox}[Домашняя работа]
  M03-H можно дать после успешной контрольной точки или после повторной
  проверки M03-R. Ошибки в заданиях 5, 6 и 7 требуют не только исправления, но и
  короткого объяснения.
\end{notebox}

\begin{teacherbox}[Ответы M03-R]
  \(8\,305\,072\); \(53\,008\); \(70\,402\,009\); \(302\,000\,070\).
  Для \(12\,005\,071\): группы \(12\), \(005\), \(071\), количества \(12\),
  \(5\), \(71\). Сборка: \(24\,006\,005\), \(7\,000\,080\). Исправления:
  \(6\,150\,302\), \(42\,007\,009\). Повторная проверка:
  \(90\,305\,008\), группа \(008\), количество \(8\), число \(12\,005\,030\).
\end{teacherbox}

\begin{teacherbox}[Ответы M03\(\star\)]
  В таблице: \(7\), \(304\), \(52\), \(18\).
  \(12005007008=12\,005\,007\,008\), группа миллионов \(005\), в числе
  \(5\) миллионов. \(4\,012\,005\,008\).

  Сравнения: \(5\,041\,008\,020<5\,041\,080\,002\), различается класс тысяч;
  \(12\,007\,500\,090<12\,070\,005\,900\), различается класс миллионов.
  Возможные числа: \(305\,001\,000\,042\) и \(305\,999\,000\,042\).
  Обратная задача: \(9\,036\,012\,000\).
\end{teacherbox}
